% Elaborado por Fabian Gomez
% Version 1.0

% =========================================================
% TEAM TET — Análisis y síntesis de requerimientos (SyRS/SRS)
% Inserción: después de s09_system_requirements
% Fuente: main.pdf (ConOps, UC, arquitectura, RF/RNF/CON, V&V)
% =========================================================

\section{Análisis y síntesis de requerimientos}\label{sec:teamtet_requirements}

\subsubsection{Introducción breve}
Este apartado consolida requerimientos del Sistema (SyRS) y de Software (SRS) verificables y trazables a ConOps y Casos de Uso, siguiendo el enfoque CONOPS $\rightarrow$ Casos de Uso $\rightarrow$ SyRS/SRS $\rightarrow$ Matrices de Verificación (I/A/D/T). El foco operacional es UC-01 (Exploración de superficie) y su ciclo diario (uplink--ejecución autónoma--downlink--idle) en sitio análogo.

% ---------- Formato estándar de requisito ----------
% (Definición centralizada en main.tex)

% =========================================================
% 1.2 SyRS — Requerimientos del sistema (SyRS-###)
% =========================================================
\subsubsection{SyRS — Requerimientos del sistema}

\paragraph{Top-level requirements}
\begin{reqblock}
  \item \ReqID{SyRS-001}
  \item \ReqShall{El sistema deberá operar como plataforma de validación de ELANav ejecutando el caso de uso UC-01 (Exploración de superficie) en un sitio análogo terrestre.}
  \item \ReqMetric{Ejecución completa del ciclo operacional (uplink $\rightarrow$ ejecución autónoma $\rightarrow$ downlink $\rightarrow$ idle) al menos 1 vez por campaña de verificación.}
  \item \ReqCond{Fase 4 (Operación nominal) en sitio análogo.}
  \item \ReqAccept{Registro de misión que evidencie: comando de inicio, ejecución autónoma y transmisión de telemetría/datos.}
  \item \ReqTrace{ConOps: Fase 4 (\S\ref{subsec:mission_phases}) y Ciclo Diario (\S\ref{subsec:daily_ops}); CU: UC-01 (\S\ref{sec:uc01}).}
\end{reqblock}

\begin{reqblock}
  \item \ReqID{SyRS-002}
  \item \ReqShall{El sistema deberá estar compuesto por los 7 subsistemas: STR, TRAC, PROP, C\&DH, EPS, COMMS y GNC.}
  \item \ReqMetric{7/7 subsistemas identificados y presentes en el baseline de arquitectura e integración.}
  \item \ReqCond{Baseline de arquitectura del sistema.}
  \item \ReqAccept{Evidencia documental del BDD/arquitectura mostrando los 7 subsistemas.}
  \item \ReqTrace{Arquitectura: \S\ref{sec:architecture}; Contexto del sistema: \S\ref{subsec:system_description}.}
\end{reqblock}

\paragraph{Functional requirements}
\begin{reqblock}
  \item \ReqID{SyRS-010}
  \item \ReqShall{El sistema deberá adquirir percepción visual mediante cámara CSI frontal (OV5647 en CAM1) para soportar la exploración autónoma.}
  \item \ReqMetric{Campo de visión horizontal según lente OV5647 instalado ($\ge$ 120° con lente M12 wide-angle; ver GNC-REQ-001).}
  \item \ReqCond{UC-01 en modo Exploración activa.}
  \item \ReqAccept{Evidencia (capturas/stream) que demuestre adquisición CSI y FOV medido conforme al lente seleccionado.}
  \item \ReqTrace{CU: UC-01 (\S\ref{sec:uc01}); Req: RF-001 / GNC-REQ-001.}
\end{reqblock}

\begin{reqblock}
  \item \ReqID{SyRS-011}
  \item \ReqShall{El sistema deberá clasificar la transitabilidad del terreno para soportar decisiones de ruta segura en UC-01.}
  \item \ReqMetric{Confianza de clasificación $\ge$ 95\% en sectores de 10$\times$10 cm.}
  \item \ReqCond{UC-01 durante Evaluar/Planificar.}
  \item \ReqAccept{Reporte de prueba con evidencia de desempeño $\ge$ 95\% bajo el criterio definido en el procedimiento de verificación.}
  \item \ReqTrace{CU: UC-01; Req: RF-002 / GNC-REQ-002.}
\end{reqblock}

\begin{reqblock}
  \item \ReqID{SyRS-012}
  \item \ReqShall{El sistema deberá detectar y evitar obstáculos frontales dentro de los umbrales de emergencia (HC-SR04 $\le$ 200\,mm o VL53L0X $\le$ 150\,mm) durante la navegación autónoma.}
  \item \ReqMetric{0 colisiones con obstáculos dentro de los umbrales de emergencia durante un recorrido de verificación.}
  \item \ReqCond{UC-01 en Desplazarse/EvitarObstáculo.}
  \item \ReqAccept{Video/telemetría + registro de eventos (\texttt{obstacle\_detected}/\texttt{cleared}) sin colisiones.}
  \item \ReqTrace{CU: UC-01 (\S\ref{sec:uc01}); Req: RF-003; Eventos: EVT-MOB001.}
\end{reqblock}

\begin{reqblock}
  \item \ReqID{SyRS-013}
  \item \ReqShall{El sistema deberá detectar deslizamiento (slip) durante la navegación y activar transición a HomingRetreat ante slip persistente.}
  \item \ReqMetric{Error de odometría $<$ 5\% respecto a ground truth (distancia real) mediante encoders (6 ruedas); transición a HomingRetreat ante slip persistente (\texttt{SLIP\_STALL\_FRAMES = 2} frames consecutivos con \texttt{stall\_mask} activo en HLC). Reemplaza criterio EKF --- ver CR-002.}
  \item \ReqCond{UC-01 en terreno con transiciones/pendientes.}
  \item \ReqAccept{Análisis comparativo (encoders vs. ground truth) con error $<$ 5\% en recorrido de referencia + verificación de transición a HomingRetreat en $\le$ 2 frames TLM ante slip inducido.}
  \item \ReqTrace{CU: UC-01 (\S\ref{sec:uc01}); Req: RF-004-R1; CR-002.}
\end{reqblock}

\begin{reqblock}
  \item \ReqID{SyRS-014}
  \item \ReqShall{El sistema deberá ejecutar una parada segura (Fault Stop) ante stall persistente de encoder, proximidad de emergencia (HC/TOF) o anomalía \texttt{OC\_FAULT}.}
  \item \ReqMetric{Activación de Fault Stop en $\le$ 500\,ms desde la detección (matches RF-005 y PROP-REQ-002).}
  \item \ReqCond{UC-01 durante Desplazarse en terreno que induzca slip o presencia de obstáculo.}
  \item \ReqAccept{Log de evento + medición temporal que muestre stall persistente (LLC \texttt{STALL\_THRESHOLD = 50} ciclos) y detención dentro de 500\,ms.}
  \item \ReqTrace{CU: UC-01 (\S\ref{sec:uc01}); Req: RF-005.}
\end{reqblock}

\begin{reqblock}
  \item \ReqID{SyRS-015}
  \item \ReqShall{El sistema deberá detener el giro de las ruedas en respuesta a una señal de Fault Stop.}
  \item \ReqMetric{Tiempo de detención $\le$ 500 ms tras la señal de Fault Stop.}
  \item \ReqCond{UC-01, cualquier estado con movimiento.}
  \item \ReqAccept{Medición con telemetría/encoder/corriente que demuestre paro $\le$ 500 ms.}
  \item \ReqTrace{Req: PROP-REQ-002 (soporte a RF-005).}
\end{reqblock}

\begin{reqblock}
  \item \ReqID{SyRS-016}
  \item \ReqShall{El sistema deberá asegurar la integridad de telemetría y comandos usando CSP como encapsulamiento sobre el enlace Rover--GCS.}
  \item \ReqMetric{100\% de paquetes TM/CMD de prueba verificados como válidos por el mecanismo de integridad (checksum/CRC) definido por CSP.}
  \item \ReqCond{UC-01 (uplink/downlink) y ciclo diario.}
  \item \ReqAccept{Captura de tráfico + verificación de integridad documentada.}
  \item \ReqTrace{CU: UC-01 (\S\ref{sec:uc01}); Req: RF-006 / COMM-REQ-001; Interfaces: \S\ref{sec:system_interfaces}.}
\end{reqblock}

\begin{reqblock}
  \item \ReqID{SyRS-017}
  \item \ReqShall{El sistema deberá reportar al C\&DH el voltaje de bus de cada banco de batería (INA3221 multi-banco) y la corriente por canal de motor (ACS712 por canal, ADC del LLC), para monitoreo distribuido de energía.}
  \item \ReqMetric{Voltaje: frecuencia $\ge$ 1\,Hz, resolución $\le$ 8\,mV (INA3221 modo bus-only). Corriente: agregada en frame TLM del LLC.}
  \item \ReqCond{Operación nominal y modo seguro.}
  \item \ReqAccept{Registro de telemetría EPS multi-banco (serie temporal) que demuestre $\ge$ 1 muestra/s y resolución de tensión $\le$ 8\,mV.}
  \item \ReqTrace{CU: UC-01 (monitorizar estado); Req: EPS-REQ-001.}
\end{reqblock}

\begin{reqblock}
  \item \ReqID{SyRS-018}
  \item \ReqShall{El sistema deberá mantener un enlace de comunicaciones en 2.4 GHz en línea de vista durante campañas de validación.}
  \item \ReqMetric{Alcance $\ge$ 50 m en línea de vista.}
  \item \ReqCond{Sitio análogo, operación nominal.}
  \item \ReqAccept{Prueba de rango con TM/CMD que evidencie enlace estable hasta $\ge$ 50 m (definición de estabilidad según procedimiento; si no existe, TBD).}
  \item \ReqTrace{Interfaces externas; Req: COMM-REQ-002.}
\end{reqblock}

\paragraph{Usability requirements}
\begin{reqblock}
  \item \ReqID{SyRS-030}
  \item \ReqShall{El sistema deberá proveer al GCS telemetría de estado suficiente para observar el modo operativo y la salud del sistema durante UC-01.}
  \item \ReqMetric{Cada mensaje de estado incluye: modo actual, timestamp, estado de enlace y telemetría de energía (V/I).}
  \item \ReqCond{UC-01 durante comunicación y ventanas de comunicación.}
  \item \ReqAccept{Captura TM que muestre campos requeridos y decodificación en GCS.}
  \item \ReqTrace{ConOps: \S\ref{subsec:daily_ops}; CU: UC-01 (\S\ref{sec:uc01}); Interfaces: \S\ref{sec:system_interfaces}.}
\end{reqblock}

\paragraph{Performance requirements}
\begin{reqblock}
  \item \ReqID{SyRS-040}
  \item \ReqShall{El sistema deberá sostener el procesamiento a bordo de agentes dentro de la latencia máxima establecida.}
  \item \ReqMetric{$\le$ 2 s por ciclo de agente.}
  \item \ReqCond{UC-01, modo Exploración activa.}
  \item \ReqAccept{Medición de latencia por ciclo en logs del sistema con estadísticos (definidos en el procedimiento) $\le$ 2 s.}
  \item \ReqTrace{CU: UC-01 (\S\ref{sec:uc01}); Req: RNF-001; Interfaces SW: \S\ref{subsec:software_interfaces}.}
\end{reqblock}

\begin{reqblock}
  \item \ReqID{SyRS-041}
  \item \ReqShall{El sistema deberá mantener la precisión de navegación dentro del error máximo permitido durante un recorrido de validación.}
  \item \ReqMetric{Error < 5\% de la distancia total recorrida.}
  \item \ReqCond{UC-01 en sitio análogo (pendientes/transiciones).}
  \item \ReqAccept{Comparación entre distancia estimada vs medición de referencia con error relativo < 5\% (método de referencia definido en el procedimiento; si no existe, TBD).}
  \item \ReqTrace{CU: UC-01; Req: RNF-003.}
\end{reqblock}

\begin{reqblock}
  \item \ReqID{SyRS-042}
  \item \ReqShall{El sistema deberá demostrar resiliencia ante contingencias operacionales mediante maniobras de recuperación/``rescate''.}
  \item \ReqMetric{Éxito $\ge$ 90\% en maniobras de rescate definidas por V\&V.}
  \item \ReqCond{UC-01 bajo contingencias (p.\,ej., atascos, pérdida de enlace, condiciones adversas).}
  \item \ReqAccept{Resultados de campaña con conteo de intentos/éxitos $\ge$ 90\% (definición de rescate según procedimiento; si no existe, TBD).}
  \item \ReqTrace{CU: UC-01; Req: RNF-002; Modos: HomingRetreat/SafeMode.}
\end{reqblock}

\paragraph{System interface requirements}
\begin{reqblock}
  \item \ReqID{SyRS-050}
  \item \ReqShall{El sistema deberá intercambiar TM/CMD con el GCS mediante UDP/IP y encapsulado CSP.}
  \item \ReqMetric{Mensajes TM/CMD decodificables como paquetes CSP transportados por UDP/IP.}
  \item \ReqCond{Uplink y downlink del ciclo diario; UC-01.}
  \item \ReqAccept{Captura de red (pcap) + decodificación CSP de mensajes de prueba.}
  \item \ReqTrace{ConOps: ciclo diario (\S\ref{subsec:daily_ops}); CU: UC-01 (\S\ref{sec:uc01}); Interfaces: \S\ref{sec:system_interfaces}.}
\end{reqblock}

\begin{reqblock}
  \item \ReqID{SyRS-051}
  \item \ReqShall{El sistema deberá permitir la inyección de latencia en el enlace Rover--GCS dentro del rango definido por la restricción CON-002.}
  \item \ReqMetric{Latencia simulada configurable entre 0 y 5\,s. \emph{Pendiente de implementación en \texttt{gcs\_drive.py} o middleware CSP.}}
  \item \ReqCond{Campañas que requieran emulación de retardo operacional.}
  \item \ReqAccept{Evidencia de configuración + medición de retardo efectivo dentro de [0, 5]\,s.}
  \item \ReqTrace{Restricción: CON-002; Interfaces: \S\ref{sec:system_interfaces}.}
\end{reqblock}

\begin{reqblock}
  \item \ReqID{SyRS-052}
  \item \ReqShall{El sistema deberá implementar las interfaces internas necesarias entre subsistemas conforme al IBD/N2 físico para soportar UC-01.}
  \item \ReqMetric{Existencia de conectividad funcional para al menos: CSI-2 (cámara), I2C (telemetría EPS), PWM/CTRL (actuación PROP) y bus de datos C\&DH$\leftrightarrow$GNC$\leftrightarrow$COMMS.}
  \item \ReqCond{Integración física del sistema.}
  \item \ReqAccept{Inspección del IBD/N2 + prueba básica (smoke test) de sensado, actuación y enlace.}
  \item \ReqTrace{Arquitectura/Interfaces: matriz N2 (Tab.~\ref{tab:n2_fisica}) y \S\ref{subsec:internal_interfaces}; CU: UC-01 (\S\ref{sec:uc01}).}
\end{reqblock}

\paragraph{Operation \& modes}
\begin{reqblock}
  \item \ReqID{SyRS-060}
  \item \ReqShall{El sistema deberá soportar los modos \texttt{Standby}, \texttt{Explore} (Exploración activa), \texttt{Avoid} (EvitarObstáculo), \texttt{Climb} (Escalada), \texttt{Retreat} (HomingRetreat), \texttt{Safe} (SafeMode) y \texttt{Fault}, conforme al \texttt{RoverState} del MSM (LLC), y transicionar entre ellos por eventos definidos. Adicionalmente, el LLC mantiene un \texttt{SafetyState} paralelo (\texttt{Normal/Warn/Limit/FaultStall}).}
  \item \ReqMetric{7 estados \texttt{RoverState} implementados; transiciones disparadas por al menos: \texttt{cmd\_start\_exploration}, \texttt{cmd\_pause\_exploration}, \texttt{cmd\_homing\_retreat}, \texttt{cmd\_climb}, \texttt{obstacle\_detected}, \texttt{sw\_fault\_critical}, \texttt{power\_anomaly\_detected}, \texttt{cmd\_end\_mission}.}
  \item \ReqCond{Operación UC-01 y contingencias.}
  \item \ReqAccept{Evidencia de ejecución de la máquina de estados (logs evento/estado) conforme a las tablas de transición.}
  \item \ReqTrace{Modelo de estados: \S\ref{subsec:uc01_states}; tabla de modos: \S\ref{subsec:op_modes}; CU: UC-01 (\S\ref{sec:uc01}).}
\end{reqblock}

\begin{reqblock}
  \item \ReqID{SyRS-061}
  \item \ReqShall{El sistema deberá ejecutar HomingRetreat al último waypoint seguro cuando se detecten condiciones adversas y la política lo permita.}
  \item \ReqMetric{Registro de selección de lastSafeWaypoint y navegación hacia él tras conditions\_adverse\_detected.}
  \item \ReqCond{UC-01; evento conditions\_\allowbreak adverse\_\allowbreak detected con guarda severity\_\allowbreak allows\_\allowbreak retreat (TBD si no está cuantificada). \textit{Impl. actual (HLC \texttt{olympus\_hlc/monitors.py::WaypointTracker.should\_retreat}):} la guarda se reduce a la conjunción \texttt{dist\_mm < RETREAT\_DIST\_MM (300\,mm) AND msm\_state != CLIMB AND not SafeMode.active}; no existe una clasificación multi-nivel de severidad. La cuantificación multi-nivel queda pendiente (ver TBD-OP-01).}
  \item \ReqAccept{Logs de transición a HomingRetreat + llegada safe\_waypoint\_reached.}
  \item \ReqTrace{Modelo de estados: \S\ref{subsec:uc01_states}; CU: UC-01 (\S\ref{sec:uc01}).}
\end{reqblock}

\paragraph{Maintenance requirements}
\begin{reqblock}
  \item \ReqID{SyRS-070}
  \item \ReqShall{El sistema deberá almacenar datos localmente para permitir recuperación y análisis posterior de campañas.}
  \item \ReqMetric{Capacidad para al menos 2 horas de telemetría y logs de misión.}
  \item \ReqCond{Operación nominal de campaña.}
  \item \ReqAccept{Verificación de almacenamiento disponible + registro continuo de 2 h sin pérdida según procedimiento.}
  \item \ReqTrace{C\&DH: CDH-REQ-002 (capacidad de almacenamiento); CONOPS: downlink/análisis.}
\end{reqblock}

\paragraph{Safety \& security requirements}
\begin{reqblock}
  \item \ReqID{SyRS-080}
  \item \ReqShall{El sistema deberá entrar a SafeMode ante falla lógica crítica o anomalía de potencia detectada.}
  \item \ReqMetric{Transición Any $\rightarrow$ SafeMode ante sw\_\allowbreak fault\_\allowbreak critical o power\_\allowbreak anomaly\_\allowbreak detected.}
  \item \ReqCond{Cualquier modo.}
  \item \ReqAccept{Logs de evento y transición; evidencia de acciones de preservación definidas en el procedimiento.}
  \item \ReqTrace{Eventos: EVT-CDH-010, EVT-EPS-005; Modelo de estados: \S\ref{subsec:uc01_states}; CU: UC-01 (\S\ref{sec:uc01}).}
\end{reqblock}

\begin{sloppypar}
\begin{reqblock}
  \item \ReqID{SyRS-081}
  \item \ReqShall{El sistema deberá implementar mecanismos de protección eléctrica distribuida ante condiciones anómalas: NTC por celda LiPo 4S, voltaje multi-banco (INA3221), corriente por motor (ACS712), fusibles 10\,A por banco y aislamiento por relés conmutados (LLC D40/D41).}
  \item \ReqMetric{Presencia y verificación funcional de las protecciones distribuidas, complementadas por el BMS Daly 4S 20A por pack (ver nota EPS en \S\ref{subsec:internal_interfaces} y CR-005).}
  \item \ReqCond{Integración del subsistema EPS.}
  \item \ReqAccept{Inspección física + verificación funcional de disparo/aislamiento conforme a procedimiento de prueba.}
  \item \ReqTrace{Arquitectura EPS: \S\ref{sec:architecture}; matriz N2: Tab.~\ref{tab:n2_fisica}.}
\end{reqblock}
\end{sloppypar}

\paragraph{Regulatory/policy requirements}
\begin{reqblock}
  \item \ReqID{SyRS-090}
  \item \ReqShall{El sistema deberá ser verificado/aceptado bajo el marco operativo lab-only / sitio análogo conforme a la restricción CON-001.}
  \item \ReqMetric{Evidencia documental de campaña en entorno análogo terrestre (lab-only/sitio análogo).}
  \item \ReqCond{Aceptación del sistema (campaña).}
  \item \ReqAccept{Acta/registro de campaña y configuración del entorno.}
  \item \ReqTrace{Restricción: CON-001; ConOps: \S\ref{subsec:operational_constraints}.}
\end{reqblock}

\begin{reqblock}
  \item \ReqID{SyRS-091}
  \item \ReqShall{El sistema deberá ser validado en un entorno que incluya pendientes de hasta 20° y obstáculos representativos para evaluar navegación y evitación.}
  \item \ReqMetric{Presencia de pendiente $\le$ 20° e inclusión de obstáculos representativos en el setup de validación.}
  \item \ReqCond{Validación de UC-01 en sitio análogo.}
  \item \ReqAccept{Evidencia del setup del sitio (medición de pendiente/obstáculos) en el reporte de campaña.}
  \item \ReqTrace{ConOps/Entorno: \S\ref{subsec:operational_environment} y \S\ref{subsec:operational_constraints}; CU: UC-01 (\S\ref{sec:uc01}).}
\end{reqblock}

\paragraph{Environmental / lifecycle / transport (si aplica)}
\begin{reqblock}
  \item \ReqID{SyRS-100}
  \item \ReqShall{El sistema deberá proteger la electrónica (C\&DH/EPS) para prevenir ingreso de partículas del sitio análogo.}
  \item \ReqMetric{Protección equivalente a IP4X (IEC 60529); método de verificación cuantitativo definido en ENV-REQ-003.}
  \item \ReqCond{Operación en sitio análogo con presencia de polvo/partículas.}
  \item \ReqAccept{Conforme a ENV-REQ-003: blow test con polvo simulado ($\le$1\,mm) durante 30\,min + inspección visual post-exposición.}
  \item \ReqTrace{STR-REQ-002 (protección de envolvente); ENV-REQ-003 (método de verificación).}
\end{reqblock}

% =========================================================
% 1.3 SRS — Requerimientos del software (SRS-###)
% =========================================================
\subsubsection{SRS — Requerimientos del software}

\paragraph{External interfaces}
\begin{reqblock}
  \item \ReqID{SRS-001}
  \item \ReqShall{El software deberá encapsular y decapsular mensajes TM/CMD usando CSP sobre UDP/IP en el enlace Rover--GCS.}
  \item \ReqMetric{Mensajes identificables como CSP y decodificables extremo a extremo.}
  \item \ReqCond{Uplink/downlink en UC-01.}
  \item \ReqAccept{Captura pcap + decodificación CSP e integridad verificada.}
  \item \ReqTrace{Interfaces SW: \S\ref{subsec:software_interfaces}; CU: UC-01 (\S\ref{sec:uc01}); Req: RF-006.}
\end{reqblock}

\begin{reqblock}
  \item \ReqID{SRS-002}
  \item \ReqShall{El software deberá permitir aplicar latencia simulada configurable entre 0 y 5\,s al intercambio TM/CMD para emular retardo operacional.}
  \item \ReqMetric{Rango configurable [0, 5]\,s. \emph{Pendiente de implementación en \texttt{gcs\_drive.py} o middleware CSP.}}
  \item \ReqCond{Campañas con CON-002 activa.}
  \item \ReqAccept{Prueba que mida el retardo impuesto dentro del rango configurado.}
  \item \ReqTrace{Restricción: CON-002; Interfaces SW: \S\ref{subsec:software_interfaces}.}
\end{reqblock}

\paragraph{Software functions}
\begin{reqblock}
  \item \ReqID{SRS-010}
  \item \ReqShall{El software deberá implementar la máquina de estados de UC-01 con estados y transiciones definidos, incluyendo el pseudoestado de history ExploracionActiva.H.}
  \item \ReqMetric{Evidencia de transición coherente con las tablas (p.\,ej., Standby $\rightarrow$ ExploracionActiva.Evaluar por cmd\_start\_exploration).}
  \item \ReqCond{Ejecución de UC-01.}
  \item \ReqAccept{Logs de estado/evento que demuestren trazas válidas para escenarios de prueba definidos.}
  \item \ReqTrace{Modelo de estados: \S\ref{subsec:uc01_states}; CU: UC-01 (\S\ref{sec:uc01}).}
\end{reqblock}

\begin{reqblock}
  \item \ReqID{SRS-011}
  \item \ReqShall{El software deberá consumir y producir los eventos definidos en el diccionario formal para UC-01.}
  \item \ReqMetric{Soporte verificable de eventos mínimos (comando/estado/falla/tiempo) según el diccionario formal.}
  \item \ReqCond{UC-01 y contingencias.}
  \item \ReqAccept{Pruebas de inyección de eventos + evidencia de respuesta/consumo y transición adecuada.}
  \item \ReqTrace{Diccionario de eventos: Tab.~\ref{tab:diccionario_eventos}; Modelo de estados: \S\ref{subsec:uc01_states}; CU: UC-01 (\S\ref{sec:uc01}).}
\end{reqblock}

\begin{reqblock}
  \item \ReqID{SRS-012}
  \item \ReqShall{El software deberá ignorar eventos no especificados para un estado dado y registrar el evento con estado, timestamp y event\_id.}
  \item \ReqMetric{Para eventos no manejados inyectados: 0 transiciones + 100\% eventos registrados con campos requeridos.}
  \item \ReqCond{Ejecución UC-01.}
  \item \ReqAccept{Logs verificables con estado, timestamp y event\_id sin cambio de estado.}
  \item \ReqTrace{Política por defecto del MSM: \S\ref{subsec:uc01_states}; CU: UC-01 (\S\ref{sec:uc01}).}
\end{reqblock}

\begin{reqblock}
  \item \ReqID{SRS-013}
  \item \ReqShall{El software deberá detectar pérdida de enlace y ejecutar la lógica de gestión/reconexión definida por los estados/eventos de GestiónEnlace.}
  \item \ReqMetric{Transición a GestiónEnlace ante link\_lost y retorno a Comunicar ante link\_restored o reconnect\_attempt\_succeeded.}
  \item \ReqCond{UC-01 durante comunicación.}
  \item \ReqAccept{Logs de transición y evidencia de reintentos conforme a guardas/políticas (si no están cuantificadas, TBD).}
  \item \ReqTrace{Modelo de estados: \S\ref{subsec:uc01_states}; Eventos: \texttt{link\_lost}/\texttt{link\_restored}; CU: UC-01 (\S\ref{sec:uc01}).}
\end{reqblock}

\begin{reqblock}
  \item \ReqID{SRS-014}
  \item \ReqShall{El software deberá detectar condición storage\_near\_full y ejecutar acciones de priorización de downlink/retención conforme a política.}
  \item \ReqMetric{Generación del evento storage\_near\_full con payload de espacio libre y priorización de transmisión cuando exista ventana de comunicación.}
  \item \ReqCond{UC-01 en GestiónDatos.}
  \item \ReqAccept{Evidencia de evento + cambio de prioridad de transmisión en presencia de comms\_window\_open.}
  \item \ReqTrace{Evento: EVT-CDH-007; Modelo de estados: GestiónDatos$\rightarrow$Comunicar; CU: UC-01.}
\end{reqblock}

\paragraph{Usability}
\begin{reqblock}
  \item \ReqID{SRS-020}
  \item \ReqShall{El software deberá formar mensajes de telemetría de estado con campos mínimos para observabilidad del operador.}
  \item \ReqMetric{4/4 campos presentes: modo, timestamp, estado de enlace y V/I.}
  \item \ReqCond{Downlink durante UC-01.}
  \item \ReqAccept{Captura y decodificación en GCS mostrando campos mínimos requeridos.}
  \item \ReqTrace{CONOPS: downlink; CU: UC-01; SyRS-030.}
\end{reqblock}

\paragraph{Performance}
\begin{reqblock}
  \item \ReqID{SRS-030}
  \item \ReqShall{El software deberá completar cada ciclo de agente en $\le$ 2 s durante UC-01.}
  \item \ReqMetric{Métrica de cumplimiento: p95 (o la definida en el procedimiento) $\le$ 2 s.}
  \item \ReqCond{Operación nominal con adquisición y planificación activas.}
  \item \ReqAccept{Perfilado/logs de tiempos por ciclo que cumplan el umbral.}
  \item \ReqTrace{Req: RNF-001; Interfaces internas SW: \S\ref{subsec:software_interfaces}; CU: UC-01 (\S\ref{sec:uc01}).}
\end{reqblock}

\paragraph{Data / database requirements}
\begin{reqblock}
  \item \ReqID{SRS-040}
  \item \ReqShall{El software deberá persistir logs y telemetría de misión en almacenamiento local para recuperación y análisis posterior.}
  \item \ReqMetric{Persistencia continua durante al menos 2 horas de operación.}
  \item \ReqCond{Campaña UC-01.}
  \item \ReqAccept{Verificación de existencia, integridad y continuidad temporal de los registros.}
  \item \ReqTrace{C\&DH: CDH-REQ-002; CONOPS: downlink/análisis; CU: UC-01.}
\end{reqblock}

\paragraph{Design constraints}
\begin{reqblock}
  \item \ReqID{SRS-050}
  % SRS-050 (alineado a Yocto Project)
  \item \ReqShall{El software \textbf{DEBERÁ} ejecutarse en un entorno de \emph{Linux embebido basado en Yocto Project} (imagen a la medida) en el subsistema C\&DH.}
  \item \ReqMetric{Verificación de imagen Yocto y servicios.}
  \item \ReqCond{Integración C\&DH y campañas.}
  \item \ReqAccept{Evidencia de la \emph{build} Yocto (manifest/lock, hashes de capas, versión de imagen) y verificación de servicios/daemons requeridos operativos en dicha imagen.}
  \item \ReqTrace{CDH-REQ-001; Suposiciones de ejecución a bordo.}
\end{reqblock}

\paragraph{Software quality attributes}
\begin{reqblock}
  \item \ReqID{SRS-060}
  \item \ReqShall{El software deberá propagar y actuar ante fallas críticas (sw\_fault\_critical) para transicionar a SafeMode y registrar evidencia.}
  \item \ReqMetric{Al inyectar sw\_fault\_critical: transición a SafeMode + registro del payload requerido.}
  \item \ReqCond{Cualquier modo.}
  \item \ReqAccept{Log de falla + transición de estado demostrable en trazas.}
  \item \ReqTrace{Evento: EVT-CDH-010; Modelo de estados: Any$\rightarrow$SafeMode; CU: UC-01.}
\end{reqblock}

\begin{reqblock}
  \item \ReqID{SRS-061}
  \item \ReqShall{El software deberá registrar eventos, transiciones y timestamps para permitir auditoría posterior de la campaña.}
  \item \ReqMetric{Cada transición registra: estado\_origen, estado\_destino, event\_id, timestamp.}
  \item \ReqCond{UC-01 y contingencias.}
  \item \ReqAccept{Revisión de logs que evidencie el esquema de auditoría en un conjunto de pruebas.}
  \item \ReqTrace{Política de logs/tabla de transiciones: \S\ref{subsec:uc01_states}; CU: UC-01 (\S\ref{sec:uc01}).}
\end{reqblock}

\subsubsection*{TBD (pendientes por definición de umbrales/políticas)}
Se mantienen como \textbf{TBD} los umbrales y políticas explícitamente no cuantificados. Estos parámetros han sido consolidados en el \textbf{TBD Register (Anexo~\ref{sec:appendices})} con su correspondiente plan de resolución, responsables y criterios de decisión, cumpliendo con la norma ISO 29148. 

Las guardas basadas en severidad y políticas de reintento se consideran TBD operativos; se definirán mediante parámetros configurables en el software (config files) documentados en el ICD de software y reflejados en los procedimientos de prueba. Su valor por defecto se definirá en la configuración de campaña conforme al plan del Anexo~\ref{sec:appendices}.
